% ============================================================
%  ChemoInteract Documentation - Part 3
%  Title: Interview Prep, Viva Q&A, Presentation & Resume Guide
%  Compile with: pdflatex ChemoInteract_Part3_Interview_and_Viva_Prep.tex
% ============================================================

\documentclass[12pt,a4paper,oneside]{article}

\usepackage[T1]{fontenc}
\usepackage[utf8]{inputenc}
\usepackage{lmodern}
\usepackage[margin=2.5cm]{geometry}
\usepackage{hyperref}
\usepackage{xcolor}
\usepackage{listings}
\usepackage{longtable}
\usepackage{booktabs}
\usepackage{array}
\usepackage{amsmath}
\usepackage{amssymb}
\usepackage{tcolorbox}
\usepackage{fancyhdr}
\usepackage{titlesec}
\usepackage{enumitem}
\usepackage{microtype}
\usepackage{setspace}

\tcbuselibrary{breakable,skins}

\hypersetup{
  colorlinks=true,
  linkcolor=blue!70!black,
  urlcolor=cyan!70!black,
  pdftitle={ChemoInteract Part 3: Interview and Viva Prep},
  pdfauthor={Ronit},
  pdfstartview=FitH,
}

% ── Colours ──────────────────────────────────────────────────
\definecolor{codegreen}{rgb}{0,0.6,0}
\definecolor{codegray}{rgb}{0.5,0.5,0.5}
\definecolor{codepurple}{rgb}{0.58,0,0.82}
\definecolor{backcolour}{rgb}{0.97,0.97,0.97}
\definecolor{deepblue}{RGB}{0,55,120}
\definecolor{chemopurple}{RGB}{91,50,160}
\definecolor{chemogreen}{RGB}{16,140,100}
\definecolor{chemoorange}{RGB}{230,100,20}
\definecolor{noteyellow}{RGB}{255,248,200}
\definecolor{tipegreen}{RGB}{210,245,215}
\definecolor{warnorange}{RGB}{255,235,210}
\definecolor{questionblue}{RGB}{210,230,255}

% ── tcolorbox styles ─────────────────────────────────────────
\tcbset{
  questionbox/.style={enhanced,breakable,colback=questionblue,
    colframe=blue!60!black,fonttitle=\bfseries,
    left=6pt,right=6pt,top=4pt,bottom=4pt},
  answerbox/.style={enhanced,breakable,colback=white,
    colframe=blue!30!black,
    fonttitle=\bfseries\color{blue!60!black},title={Model Answer},
    left=6pt,right=6pt},
  vivabox/.style={enhanced,breakable,colback=questionblue!40,
    colframe=deepblue,fonttitle=\bfseries,title={Viva Question},
    left=6pt,right=6pt},
  cheatbox/.style={enhanced,breakable,colback=chemopurple!8,
    colframe=chemopurple,fonttitle=\bfseries\color{white},
    coltitle=white,attach boxed title to top left,
    boxed title style={colback=chemopurple,rounded corners},
    left=6pt,right=6pt},
}

% ── Header / Footer ──────────────────────────────────────────
\pagestyle{fancy}
\fancyhf{}
\fancyhead[L]{\small\textcolor{chemopurple}{\textbf{ChemoInteract --- Part 3}}}
\fancyhead[R]{\small\textcolor{gray}{Interview, Viva \& Presentation Guide}}
\fancyfoot[C]{\thepage}
\renewcommand{\headrulewidth}{0.5pt}

% ── Section styles ───────────────────────────────────────────
\titleformat{\section}{\Large\bfseries\color{deepblue}}{\thesection}{1em}{}
\titleformat{\subsection}{\large\bfseries\color{chemopurple}}{\thesubsection}{1em}{}
\titleformat{\subsubsection}{\normalsize\bfseries\color{chemogreen}}{\thesubsubsection}{1em}{}

\setstretch{1.2}
\setlength{\parskip}{6pt}

\begin{document}

% ── Title Block ──────────────────────────────────────────────
\begin{center}
{\Huge\bfseries\color{deepblue} ChemoInteract}\\[0.3cm]
{\Large\bfseries\color{chemopurple} Complete Documentation --- Part 3 of 3}\\[0.2cm]
{\large\textbf{Interview Q\&A, College Viva Preparation, Presentation Scripts, Resume & Roadmap}}\\[0.4cm]
\begin{tcolorbox}[enhanced,width=14cm,colback=deepblue!5,colframe=deepblue,boxrule=1pt,arc=5pt]
\centering
\begin{tabular}{llll}
\textbf{Author:} & Ronit & \textbf{Version:} & 0.1.0 \\
\textbf{Framework:} & PyTorch 2.0+ & \textbf{Focus:} & Interview \& Viva Prep \\
\end{tabular}
\end{tcolorbox}
\end{center}

\tableofcontents
\newpage

% ================================================================
\section{Interview Preparation Guide}\label{sec:interview}

\subsection{HR \& Project Overview Questions}

\begin{tcolorbox}[questionbox,title={Q1: Tell me about your final year project.}]
\end{tcolorbox}
\begin{tcolorbox}[answerbox]
``I developed ChemoInteract, a modular deep learning Python library for predicting drug-drug interactions. Adverse drug reactions cause thousands of preventable hospitalizations annually in polypharmacy patients. ChemoInteract unifies 11 published deep learning architectures --- ranging from simple FC networks to Graph Neural Networks and cross-drug Attention models --- under a single plug-and-play training pipeline. A researcher can benchmark any model across 5 auto-downloaded datasets with a single line of code.''
\end{tcolorbox}

\begin{tcolorbox}[questionbox,title={Q2: What was your biggest technical challenge?}]
\end{tcolorbox}
\begin{tcolorbox}[answerbox]
``Making 11 models with completely different input requirements work through one unified pipeline. Some take fingerprint vectors; others take molecular graphs. Some need cell line context; others do not. I solved this by implementing the \textbf{Template Method} pattern via an abstract \texttt{unpack()} method on the \texttt{Model} base class. Each model declares what it extracts from the \texttt{DrugPairBatch}, keeping the pipeline completely decoupled from model internals.''
\end{tcolorbox}

\subsection{Technical Questions & Model Answers}

\begin{tcolorbox}[questionbox,title={Q3: Why use ROC-AUC instead of Accuracy?}]
\end{tcolorbox}
\begin{tcolorbox}[answerbox]
``Drug interaction datasets are heavily class-imbalanced (far more non-interacting pairs than interacting ones). A naive model predicting all zeros gets high accuracy but is clinically useless. ROC-AUC measures the probability that a randomly chosen positive pair receives a higher prediction score than a randomly chosen negative pair, making it threshold-independent and robust to class imbalance.''
\end{tcolorbox}

\begin{tcolorbox}[questionbox,title={Q4: Explain \texttt{@lru\_cache} in your data loader.}]
\end{tcolorbox}
\begin{tcolorbox}[answerbox]
``\texttt{@lru\_cache(maxsize=1)} memoises dataset loading methods. The first call downloads and parses data from GitHub/disk; subsequent calls return the cached in-memory objects instantly, avoiding redundant network I/O during repeated training runs.''
\end{tcolorbox}

% ================================================================
\section{College Viva Questions and Answers}\label{sec:viva}

\begin{tcolorbox}[vivabox]
V1: What does DDI stand for?
\end{tcolorbox}
Drug-Drug Interaction: when the pharmacological effect of one drug is altered by another drug co-administered in the body.

\begin{tcolorbox}[vivabox]
V2: How many deep learning models are in ChemoInteract?
\end{tcolorbox}
11 models: DeepSynergy, DeepDDI, MatchMaker (FC models); DeepDDS, DeepDrug, CASTER, GCNBMP, EPGCN-DS, MR-GNN (GNN models); MHCADDI, SSI-DDI (Attention models).

\begin{tcolorbox}[vivabox]
V3: What is a SMILES string?
\end{tcolorbox}
Simplified Molecular Input Line Entry System: an ASCII text representation of chemical structures (e.g. Aspirin is \texttt{CC(=O)Oc1ccccc1C(=O)O}). TorchDrug parses SMILES strings into molecular graphs.

\begin{tcolorbox}[vivabox]
V4: What is a Morgan Fingerprint?
\end{tcolorbox}
A 256-bit binary vector encoding the presence of circular chemical substructures up to radius 2 bonds around each atom.

\begin{tcolorbox}[vivabox]
V5: What does \texttt{model.eval()} do?
\end{tcolorbox}
Disables Dropout (all neurons stay active) and locks Batch Normalization to use stored running statistics during evaluation.

\begin{tcolorbox}[vivabox]
V6: What is the Closed-World Assumption (CWA)?
\end{tcolorbox}
Assuming that any unrecorded drug pair in database records is non-interacting (label = 0.0) during negative sampling.

% ================================================================
\section{Presentation Scripts}\label{sec:presentation}

\subsection{2-Minute Elevator Pitch}
``I built ChemoInteract, a PyTorch library using deep learning to predict dangerous drug-drug interactions. Adverse drug reactions cause thousands of hospitalizations annually.

My primary contribution is unifying 11 published deep learning architectures --- from basic neural networks to GNNs with cross-drug co-attention --- into a single framework. Any model can be trained on any of 5 benchmark datasets with a single line of code.

The library automatically handles GPU acceleration, dataset downloading, and metric evaluation, achieving ROC-AUC scores up to 0.87 on benchmark datasets.''

\subsection{30-Second Pitch}
``ChemoInteract is an open-source PyTorch library that uses Graph Neural Networks and AI to predict drug interactions and synergies. It unifies 11 published deep learning models across 5 datasets in a single line of code, enabling rapid computational drug discovery.''

% ================================================================
\section{Resume Preparation}\label{sec:resume}

\subsection{Resume Bullet Points}
\begin{itemize}
  \item \textbf{ChemoInteract Framework:} Designed and developed a modular PyTorch library unifying \textbf{11 published deep learning models} (GNNs, MLPs, Co-Attention) for drug interaction, polypharmacy side-effect, and synergy prediction.
  \item \textbf{Modular Pipeline:} Implemented Template Method pattern via abstract \texttt{unpack()} method, reducing experiment training code from 100+ lines of boilerplate to a 1-line \texttt{pipeline()} call.
  \item \textbf{Compatibility Layer:} Engineered GPU compatibility patches for TorchDrug \texttt{PackedGraph} objects, resolving tensor transfer issues across PyTorch versions.
  \item \textbf{Benchmarking \& Testing:} Integrated 5 auto-cached datasets with \texttt{@lru\_cache} and built a 4-file \texttt{pytest} suite verifying forward passes and data structures.
\end{itemize}

\subsection{ATS Keywords}
PyTorch, Deep Learning, Graph Neural Networks (GNN), GCN, Drug-Drug Interaction (DDI), Molecular Representation Learning, RDKit, SMILES, Morgan Fingerprints, Co-Attention, Binary Cross-Entropy, ROC-AUC, scikit-learn, Python, TorchDrug, Unit Testing, Modular Software Architecture.

% ================================================================
\section{Future Roadmap & Quick Cheat Sheet}\label{sec:cheatsheet}

\begin{tcolorbox}[cheatbox,title={ChemoInteract Quick Reference Cheat Sheet}]
\textbf{Library Entrypoint:} \texttt{from chemointeract import pipeline} \\
\textbf{Execution:} \texttt{results = pipeline(dataset="drugcombdb", model="deepsynergy", epochs=100)} \\
\textbf{Core Design Patterns:} Template Method (\texttt{unpack()}), Resolver Pattern (\texttt{class-resolver}), Caching (\texttt{@lru\_cache}), Monkey-patch Shim (\texttt{compat.py}). \\
\textbf{Key Datasets:} DrugCombDB, DrugComb, TwoSides, DrugbankDDI, OncoPolyPharmacology. \\
\textbf{Top Models:} DeepSynergy (FC baseline), DeepDDS (GCN), MHCADDI (Cross-drug Attention), GCNBMP (FFT Circular Correlation).
\end{tcolorbox}

\end{document}
